\paragraph{Búsqueda de videojuegos}
Desde la pantalla de inicio de la aplicación, se pueden buscar los videojuego por título, género, desarrolladora y/o año de lanzamiento(RF-5.1).
En el frontend, la pantalla principal presenta dos componentes. Por un lado, el componente VideogameFilter que presenta un formulario reactivo de búsqueda.
En dicho componente, cuando el usuario rellena el formulario 
y pulsa el botón de búsqueda, se genera un evento mediante EventEmitter, que llama al servicio SearchService, el cual realiza 
la petición con los parámetros de búsqueda al endpoint 'videojuegos/'. Una vez que se obtiene la respuesta del servidor y se mapea a objetos comprensibles, 
entra en juego el componente VideogameList, que muestra los videojuegos encontrados en forma de tarjetas. 

En el servidor, el viewset se configura con DjangoFilterBackend y SearchFilter, lo que permite realizar búsquedas que coincidan con el título o 
desarrolladora del videojuego. En la figura \ref{FIG:SEARCH} se muestran los resultados de coincidencias con la desarrolladora 'Nintendo'.

\begin{figure}[Coincidencias en desarrolladora]{FIG:SEARCH}{Salida de la búsqueda de videojuegos por desarrolladora escribiendo el nombre incompleto}
    \image{13cm}{}{search}
\end{figure}
Si no se recibe ningún parámetro de búsqueda, se devuelven todos los videojuegos disponibles. La figura 
\ref{FIG:DETALLE} muestra la interacción de los componentes.

\begin{figure}[Componentes implicados para la pantalla de detalle de videojuego]{FIG:DETALLE}{ En este diagrama se puede ver el proceso completo de la pantalla de detalle de videojuego}
    \image{13cm}{}{detalleVideojuego}
\end{figure}

\paragraph{Detalle de videojuego}
Dicho sistema también se presenta en la vista de detalle de cada videojuego(RF-5.2). Dicha pantalla viene enlazada con el subsistema de valoración 
ya que en ella se pueden ver las valoraciones de los usuarios y realizarlas y con el subsistema de listas ya que se pueden añadir videojuegos 
a las listas personalizadas. Desde el fronted, el componente VideogameDetailComponent muestra la información del videojuego mediante el objeto 
Videogame (mapeado anteriormente una vez obtenida la respuesta del servidor). Dicha vista realiza una petición al servicio SearchService al 
endpoint 'videojuegos/{id}', donde {id} es el identificador del videojuego.

Por el lado del backend, simplemente se obtiene el videojuego con el id indicado en la petición y se devuelve como respuesta.

El flujo de interacción se observa en la figura \ref{FIG:BUS}.

\begin{figure}[Componentes implicados para la pantalla de inicio]{FIG:BUS}{ En este diagrama se puede ver el proceso completo de la búsqueda de videojuegos}
    \image{13cm}{}{busqueda}
\end{figure}
